\documentclass[conference]{IEEEtran}
\IEEEoverridecommandlockouts

% ── Packages ──────────────────────────────────────────────────────────────────
\usepackage{cite}
\usepackage{amsmath,amssymb,amsfonts}
\usepackage{graphicx}
\usepackage{textcomp}
\usepackage{xcolor}
\usepackage{booktabs}
\usepackage{array}
\usepackage{balance}

% ── Title ─────────────────────────────────────────────────────────────────────
\begin{document}

\title{Performance Analysis of Grover's Search Algorithm\\
Under Realistic Noise Models on NISQ Simulators}

\author{
  \IEEEauthorblockN{Saim Haider}
  \IEEEauthorblockA{
    \textit{Department of Computer Science}\\
    \textit{FAST-NUCES, Pakistan}\\
    22P-9244\\
    p229244@pwr.nu.edu.pk
  }
  \and
  \IEEEauthorblockN{Muhammad Abdullah}
  \IEEEauthorblockA{
    \textit{Department of Computer Science}\\
    \textit{FAST-NUCES, Pakistan}\\
    22P-9371\\
    p229371@pwr.nu.edu.pk
  }
  \and
  \IEEEauthorblockN{Abdullah}
  \IEEEauthorblockA{
    \textit{Department of Computer Science}\\
    \textit{FAST-NUCES, Pakistan}\\
    22P-9358\\
    p229358@pwr.nu.edu.pk
  }
}

\maketitle

% ── Abstract ──────────────────────────────────────────────────────────────────
\begin{abstract}
Grover's search algorithm provides a provable quadratic speedup over classical
unstructured search, making it one of the most relevant quantum algorithms for
near-term deployment. However, current quantum hardware operates in the Noisy
Intermediate-Scale Quantum (NISQ) regime, where decoherence and gate errors
severely degrade algorithm performance. This research presents a simulation-based
comparative study of Grover's algorithm under three realistic noise models:
depolarizing noise, bit-flip errors, and phase-flip errors. Using Python and the
Qiskit Aer simulator, Grover's search circuit is implemented for search spaces of
$n = 2, 3, 4, 5$ qubits. Noise is injected at parameterized error rates swept
across the range $p \in [0.001, 0.1]$ and success probability is measured at each
configuration. The primary objective is to identify the critical noise threshold
$p^{*}$ at which Grover's algorithm loses its quantum advantage over classical
brute-force search, and to determine how this threshold scales with the number of
qubits. Results include empirical noise threshold curves for each noise type,
crossover analysis comparing Grover's performance against the classical baseline
$\frac{1}{N}$, and a reproducible Qiskit codebase. This work directly addresses
a gap in the existing literature by systematically identifying quantum advantage
loss thresholds across multiple noise types using modern quantum simulation tools.
\end{abstract}

\begin{IEEEkeywords}
Grover's algorithm, NISQ computing, noise models, quantum advantage
\end{IEEEkeywords}

% ── I. Introduction ───────────────────────────────────────────────────────────
\section{Introduction}

Quantum computing leverages the principles of superposition, entanglement, and
quantum interference to solve certain computational problems significantly faster
than classical computers. Among the most practically relevant quantum algorithms
is Grover's search algorithm~\cite{grover1996}, which achieves
$\mathcal{O}(\sqrt{N})$ query complexity for searching an unstructured database
of $N$ entries, compared to the classical $\mathcal{O}(N)$ average case. This
quadratic speedup has wide applications in optimization, cryptanalysis, and
database search, making Grover's algorithm a standard benchmark for evaluating
near-term quantum hardware.

Despite its theoretical promise, quantum algorithms face a fundamental challenge
in practice: noise. Physical qubits are susceptible to environmental decoherence,
gate imperfections, and measurement errors. These noise sources collectively
corrupt the quantum interference patterns that Grover's algorithm depends upon to
amplify the probability of measuring the correct answer. Present-day quantum
processors operate in the Noisy Intermediate-Scale Quantum (NISQ) regime,
characterized by gate error rates on the order of $10^{-3}$ to $10^{-2}$ and
coherence times in the microsecond to millisecond range. Under these conditions,
even moderately deep circuits such as Grover's can suffer severe degradation in
success probability.

While previous research has studied noise effects on Grover's algorithm, no work
has systematically combined modern Qiskit-based circuit simulation with explicit
identification of the quantum advantage loss threshold $p^{*}$ across multiple
isolated noise channels. This paper analyzes a systematic empirical study of
Grover's algorithm under realistic noise conditions using the Qiskit Aer
simulator. By sweeping noise error rates across a controlled range and measuring
success probability for each configuration, this work aims to identify the precise
noise threshold beyond which Grover's algorithm loses its quantum advantage over
classical search.

\subsection{Contributions}

This work extends the current literature by making the following specific
contributions:

\begin{enumerate}
  \item Implementation of a comprehensive Qiskit-based simulation framework for
        Grover's algorithm across $n = 2, 3, 4, 5$ qubit configurations with
        systematic noise injection.

  \item Empirical identification of the quantum advantage threshold $p^{*}$ for
        three distinct noise models through controlled error rate sweeps from
        $p = 0.001$ to $p = 0.1$.

  \item Explicit crossover analysis comparing Grover's noisy performance against
        the classical random search baseline $\frac{1}{2^{n}}$.

  \item Development of a reproducible open-source Qiskit codebase suitable for
        extension to noise mitigation and error correction research.
\end{enumerate}

The remainder of this paper is organized as follows: Section~II reviews related
work and identifies gaps in existing literature. Section~III concludes with a
summary of expected outcomes and future research directions.

% ── II. Literature Review ─────────────────────────────────────────────────────
\section{Literature Review}

This section reviews prior work on noise effects in Grover's algorithm and
identifies the specific research gap this work addresses.

\subsection{Ijaz and Faryad (2023) -- Base Paper}

Ijaz and Faryad~\cite{ijaz2023} conducted a comprehensive noise analysis of
Grover's algorithm using Qiskit Aer for $n = 3, 4, 5$ qubits with 10{,}000 shots
per configuration. Four noise channels were studied: depolarizing, bit-flip (X),
phase-flip (Z), and bit-phase flip (Y). The key finding was that success
probability decays exponentially with error probability $p$ but only linearly with
qubit count $n$. The bit-phase flip channel caused the most severe degradation
across all configurations.

\textit{Strength:} Real Qiskit Aer simulation with four noise types providing
empirical confirmation of exponential decay relationship between success
probability and error rate.

\textit{Limitation:} Does not identify the critical threshold $p^{*}$ where
quantum advantage over classical search is lost. The study focuses on comparing
standard versus QSVT implementations rather than explicit threshold analysis.

\subsection{Salas (2008)}

Salas~\cite{salas2008} performed the earliest systematic numerical study of noise
in Grover's algorithm using Fortran-based Monte Carlo simulation for $n = 2$ to
$7$ qubits under a depolarizing channel model. The study derived an exponential
damping law for success probability and an allowed-error law:
$\varepsilon_{\text{th}}(N) \sim N^{-1.1}$, where $N = 2^{n}$. An absolute
threshold of $\varepsilon > 0.043$ was identified above which the algorithm fails
for any qubit count.

\textit{Strength:} First quantitative noise threshold derived for Grover's
algorithm providing foundational analytical results that remain cited in modern
literature.

\textit{Limitation:} Uses an outdated Fortran framework predating Qiskit. Studies
only depolarizing noise without explicit comparison against classical search
advantage boundary.

\subsection{Dowarah et al.\ (2025)}

Dowarah et al.~\cite{dowarah2025} applied Random Matrix Theory (RMT) to model the
noisy Grover operator as a Floquet unitary and derived two analytical noise
thresholds. The computing transition at $\delta_{c,\text{comp}} \sim
\mathcal{O}(2^{-L/2})$ marks where quantum advantage is lost, and the ergodicity
transition at $\delta_{c,\text{gap}} \sim \mathcal{O}(L^{-1})$ marks where circuit
structure collapses entirely.

\textit{Strength:} State-of-the-art analytical framework identifying two distinct
thresholds with exact scaling laws based on rigorous theoretical foundations.

\textit{Limitation:} Entirely theoretical work restricted to systematic coherent
noise. No Qiskit simulation or empirical validation provided to confirm predictions
on actual quantum circuits.

\subsection{Srivastava et al.\ (2025)}

Srivastava et al.~\cite{srivastava2025} proposed using quantum switches exploiting
indefinite causal order to mitigate depolarizing noise in Grover's algorithm. Two
frameworks were developed, with the second framework (delayed measurement)
providing significantly greater noise resilience by preserving correlations between
the input state and quantum switch register throughout all iterations.

\textit{Strength:} Novel mitigation approach demonstrating improved success
probability analytically. Framework is directly applicable to NISQ devices for
practical noise reduction.

\textit{Limitation:} Entirely analytical with no Qiskit simulation to validate
performance improvements. Studies only depolarizing noise and does not identify
the threshold $p^{*}$ for standard Grover circuits without mitigation.

\subsection{Comparative Analysis and Research Gap}

Table~\ref{tab:comparison} summarizes the strengths and limitations of reviewed
literature. The comparative analysis reveals that no existing study combines
(1) a modern Qiskit simulation framework, (2) multiple isolated noise types,
(3) systematic error rate sweeps, and (4) explicit identification of the quantum
advantage loss threshold.

\begin{table}[htbp]
  \centering
  \caption{Comparative Analysis of Related Work}
  \label{tab:comparison}
  \resizebox{\columnwidth}{!}{%
  \begin{tabular}{p{2.0cm} p{1.8cm} p{2.2cm} p{2.2cm}}
    \toprule
    \textbf{Paper} & \textbf{Method} & \textbf{Strength} & \textbf{Limitation} \\
    \midrule
    Ijaz \& Faryad (2023)~\cite{ijaz2023}
      & Qiskit Aer simulation
      & 4 noise types; exponential decay confirmed
      & No $p^{*}$ threshold; QSVT focus \\[4pt]

    Salas (2008)~\cite{salas2008}
      & Monte Carlo, Fortran
      & First threshold law derived
      & Outdated tools; depolarizing only \\[4pt]

    Dowarah et al.\ (2025)~\cite{dowarah2025}
      & Random Matrix Theory
      & Two-threshold analytical model
      & No simulation; coherent noise only \\[4pt]

    Srivastava et al.\ (2025)~\cite{srivastava2025}
      & Analytical, quantum switch
      & ICO mitigation; improved success
      & No Qiskit; depolarizing only \\
    \bottomrule
  \end{tabular}}
\end{table}

This work fills the identified gap by implementing a systematic Qiskit-based
empirical study that explicitly locates the quantum advantage threshold $p^{*}$
across three noise types through controlled error sweeps, providing concrete
benchmarks for NISQ hardware calibration.

% ── III. Conclusion ───────────────────────────────────────────────────────────
\section{Conclusion}

This research addresses a critical gap in quantum computing literature by
systematically identifying the noise threshold at which Grover's search algorithm
loses its quantum computational advantage over classical brute-force search. By
combining modern Qiskit-based simulation with comprehensive noise sweeps across
depolarizing, bit-flip, and phase-flip channels, this work will deliver empirical
threshold curves that are directly actionable for NISQ hardware development and
calibration strategies.

The expected outcomes include: (1) empirical noise threshold values $p^{*}$ for
each noise type and qubit configuration, (2) scaling analysis showing how
thresholds evolve with increasing system size, (3) crossover plots explicitly
marking the quantum advantage boundary, and (4) a reproducible open-source Qiskit
codebase suitable for extension to error mitigation research.

Compared to Ijaz and Faryad (2023), this work will produce the explicit quantum
advantage threshold map their study does not provide. Compared to Salas (2008),
results will be validated on modern Qiskit circuit simulators with three noise
types rather than one. The threshold values obtained will serve as direct empirical
inputs for noise mitigation strategies such as the quantum switch framework of
Srivastava et al.\ (2025), and will provide the circuit-level validation that the
analytical RMT framework of Dowarah et al.\ (2025) presupposes.

This work establishes a foundation for future research in quantum algorithm
robustness under realistic noise conditions, providing concrete benchmarks for
near-term quantum hardware capabilities and limitations.

% ── References ────────────────────────────────────────────────────────────────
\balance
\begin{thebibliography}{00}

\bibitem{grover1996}
L.~K. Grover, ``A fast quantum mechanical algorithm for database search,''
in \textit{Proc. 28th Annual ACM Symposium on Theory of Computing (STOC)},
1996, pp. 212--219.

\bibitem{ijaz2023}
M.~A. Ijaz and M.~Faryad, ``Noise analysis of Grover and phase estimation
algorithms implemented as quantum singular value transformations for a small
number of noisy qubits,'' \textit{Scientific Reports}, vol.~13, p.~20144, 2023.

\bibitem{salas2008}
P.~J. Salas, ``Noise effect on Grover algorithm,''
\textit{The European Physical Journal D}, vol.~46, pp.~365--381, 2008.

\bibitem{dowarah2025}
S.~Dowarah, C.~Zhang, V.~Khemani, and M.~Kolodrubetz,
``Phases and phase transition in Grover's algorithm with systematic noise,''
\textit{arXiv preprint arXiv:2406.10344v2}, 2025.

\bibitem{srivastava2025}
S.~Srivastava, A.~K. Pati, I.~Chakrabarty, and S.~Bhattacharya,
``Using quantum switches to mitigate noise in Grover's search algorithm,''
\textit{arXiv preprint arXiv:2401.05866v2}, 2025.

\end{thebibliography}

\end{document}